\section*{Chapitre \ref{ch:b1:alcohol-activation} --- Les alcools~: activer le groupe OH}

\begin{solution}{exo:b1:alcohol-activation:1}
\ce{2CH3CH2OH + 2Na -> 2CH3CH2O- + 2Na+ + H2}~;
\ce{CH3CH2OH + NaH -> CH3CH2O- + Na+ + H2}~;
\ce{CH3CH2OH + OH- <=> CH3CH2O- + H2O}, $K = 10^{14.0 - 15.9} = 10^{-1.9}$~:
réaction partielle.
\end{solution}

\begin{solution}{exo:b1:alcohol-activation:2}
1-Méthoxypropane~; méthoxyéthane~; éthoxybenzène (le phénolate, base plus
faible qu'un \omterm{def:b1:alcohol-activation:alkoxide}{alcoolate}, reste un bon \omterm{def:b1:electronic-effects:nucleophile}{nucléophile}).
\end{solution}

\begin{solution}{exo:b1:alcohol-activation:3}
La tosylation dans la pyridine donne \ce{CH3CH2CH2OTs}~; puis
\[ \ce{CH3CH2CH2OTs + I- -> CH3CH2CH2I + TsO-}. \] En réaction directe,
l'iodure devrait chasser \ce{OH-}~: aucune réaction. Le \omterm{def:b1:alcohol-activation:sulfonate}{tosylate} possède
un excellent \omterm{def:b1:electronic-effects:nucleophile}{groupe partant}, et le milieu n'est ni acide ni fortement
basique.
\end{solution}

\begin{solution}{exo:b1:alcohol-activation:4}
Butan-2-ol~: but-1-ène, \cip{Z}- et \cip{E}-but-2-ène~; majoritaire~:
\cip{E}-but-2-ène. 2-Méthylpentan-2-ol~: 2-méthylpent-2-ène (majoritaire,
trisubstitué) et 2-méthylpent-1-ène.
\end{solution}

\begin{solution}{exo:b1:alcohol-activation:5}
\ce{(CH3)2CHCH2-O-CH2CH3}~: le 2-méthylpropan-1-olate de sodium avec le
bromoéthane, ou l'éthanolate de sodium avec le 1-bromo-2-méthylpropane.
Les deux halogénures sont primaires, mais le 1-bromo-2-méthylpropane est
ramifié à côté du carbone qui réagit, ce qui ralentit la SN2 et laisse la
\omterm{def:b1:predominant-reaction:strong-acid}{base forte} éliminer (2-méthylpropène). Le premier choix est le meilleur.
\end{solution}

\begin{solution}{exo:b1:alcohol-activation:6}
Le \omterm{def:b1:alcohol-activation:sulfonate}{tosylate} garde la \omterm{def:b1:stereochemistry-in-depth:isomers}{configuration} \cip{R}~; la SN2 avec l'éthanolate
l'inverse~: \cip{S}-2-éthoxyoctane (OEt est prioritaire comme l'était
OTs). Avec l'acide sulfurique à chaud~: élimination en octènes
(surtout l'\cip{E}-oct-2-ène), par un \omterm{def:b1:electronic-effects:intermediates}{carbocation} secondaire, avec perte
de la stéréochimie.
\end{solution}

\begin{solution}{exo:b1:alcohol-activation:7}
Protonation du OH~; perte d'eau conduisant au \omterm{def:b1:electronic-effects:intermediates}{carbocation} tertiaire
(étape lente)~; capture par \ce{Cl-}~: 2-chloro-2-méthylpropane. La
réaction est rapide car le \omterm{def:b1:electronic-effects:intermediates}{carbocation} tertiaire est stable, et le
chlorure, insoluble dans l'acide, se sépare.
\end{solution}

\begin{solution}{exo:b1:alcohol-activation:8}
\ce{CH3CH2OH + H+ -> CH3CH2OH2+}~; une seconde molécule d'éthanol attaque
le carbone par l'arrière tandis que l'eau part (SN2)~:
\ce{(CH3CH2)2OH+}, qui perd un proton~: éthoxyéthane. Réaction
concurrente~: une base (l'éthanol, \ce{HSO4-}) arrache un proton en
$\beta$ de \ce{CH3CH2OH2+} pendant que l'eau part (E2)~: éthène.
\end{solution}

\begin{solution}{exo:b1:alcohol-activation:9}
Le 2-méthylpropan-2-ol (tertiaire), qui donne le 2-chloro-2-méthylpropane, insoluble.
\end{solution}

\begin{solution}{exo:b1:alcohol-activation:10}
La protonation et la perte d'eau donnent un \omterm{def:b1:electronic-effects:intermediates}{carbocation} secondaire voisin
d'un carbone quaternaire. Un groupe méthyle de ce carbone migre avec son
\omterm{def:b1:lewis-resonance-vsepr:lewis}{doublet liant} vers le carbone cationique~: la charge positive se trouve
alors sur un \omterm{def:b1:electronic-effects:carbon-class}{carbone tertiaire}, plus stable. La perte d'un proton d'un
carbone voisin donne le 2,3-diméthylbut-2-ène, tétrasubstitué, l'alcène
le plus stable.
\end{solution}

\begin{solution}{exo:b1:alcohol-activation:11}
À partir du \cip{R}-butan-2-ol~: \omterm{def:b1:alcohol-activation:sulfonate}{tosylate} (\cip{R}), puis méthanolate de
sodium (SN2, inversion)~: \cip{S}-2-méthoxybutane. À partir du
\cip{S}-butan-2-ol~: former l'\omterm{def:b1:alcohol-activation:alkoxide}{alcoolate} (NaH) et ajouter l'iodométhane~;
la liaison C--O du carbone stéréogène n'est pas touchée~: on obtient de
nouveau le \cip{S}-2-méthoxybutane.
\end{solution}

\begin{solution}{exo:b1:alcohol-activation:12}
$Q = [\text{éther}][\ce{H2O}]/[\ce{EtOH}]^2$. Éliminer l'eau à mesure
qu'elle se forme maintient $[\ce{H2O}]$ et $Q$ faibles~: $Q < K$, et la
réaction continue de former l'éther jusqu'à consommation de l'alcool.
\end{solution}

\begin{solution}{pb:b1:alcohol-activation:1}
\textbf{1.} Le \textit{tert}-butanolate de sodium avec un halogénure de
méthyle~; le méthanolate de sodium avec un halogénure de
\textit{tert}-butyle.
\textbf{2.} Le second~: le méthanolate, \omterm{def:b1:predominant-reaction:strong-acid}{base forte}, provoque
l'élimination sur l'halogénure tertiaire,
\ce{(CH3)3CBr + CH3O- -> (CH3)2C=CH2 + CH3OH + Br-}.
\textbf{3.} \ce{(CH3)3CO- + CH3I -> (CH3)3COCH3 + I-}~: l'oxygène de
l'\omterm{def:b1:alcohol-activation:alkoxide}{alcoolate} attaque le carbone du méthyle par l'arrière pendant que
l'iodure part (SN2).
\textbf{4.} Avec le sodium (ou l'hydrure de sodium)~: \ce{2(CH3)3COH + 2Na ->
2(CH3)3CO- + 2Na+ + H2}. L'hydroxyde est une base trop faible~: $K = 10^{14.0 -
19.2} = 10^{-5.2}$.
\textbf{5.} L'eau protonerait l'\omterm{def:b1:alcohol-activation:alkoxide}{alcoolate} et attaquerait l'halogénure.
\textbf{6.} Non~: il n'a pas de \omterm{def:b1:stereochemistry-in-depth:stereocentre}{centre stéréogène}.
\textbf{7.} \cip{R}-butan-2-ol + TsCl (pyridine) $\to$ \omterm{def:b1:alcohol-activation:sulfonate}{tosylate} de
\cip{R}-butan-2-yle~: rétention.
\textbf{8.} \cip{S}-2-méthoxybutane, par inversion SN2.
\textbf{9.} Le \cip{S}-butan-2-ol.
\textbf{10.} \cip{R}-2-méthoxybutane~: la liaison C--O du carbone
stéréogène est conservée~; seule la liaison O--C(méthyle) se forme.
\textbf{11.} L'E2 (le méthanolate est une \omterm{def:b1:predominant-reaction:strong-acid}{base forte} et le carbone est
secondaire), qui donne des butènes~; elle est limitée par la petite
taille de la base, l'excellent \omterm{def:b1:electronic-effects:nucleophile}{groupe partant} et une température modérée.
\textbf{12.} L'alcool protoné s'ionise en un \omterm{def:b1:electronic-effects:intermediates}{carbocation} secondaire plan,
capturé sur ses deux faces par le méthanol~: éther racémique, accompagné
de butènes.
\textbf{13.} Protonation du OH, perte d'eau (étape lente) conduisant au
\omterm{def:b1:electronic-effects:intermediates}{carbocation} tertiaire, perte d'un proton en $\beta$.
\textbf{14.} Le 2-méthylbut-2-ène (majoritaire, trisubstitué) et le
2-méthylbut-1-ène.
\textbf{15.} Un \omterm{def:b1:electronic-effects:intermediates}{carbocation} tertiaire est encombré et perd un proton bien
plus facilement qu'une autre molécule d'alcool ne peut l'attaquer.
\textbf{16.} Pour éliminer un produit~: $Q$ restant faible, la réaction
se poursuit, et l'alcène ne reste pas au contact de l'acide.
\textbf{17.} Une première barrière élevée (perte d'eau, \omterm{def:b1:elementary-steps:rds}{étape
cinétiquement déterminante}) jusqu'au \omterm{def:b1:electronic-effects:intermediates}{carbocation}, puis une petite
barrière jusqu'à l'alcène.
\textbf{18.} Non~: un \omterm{def:b1:electronic-effects:intermediates}{carbocation} primaire est trop instable~; les
alcools primaires éliminent par E2 sur l'alcool protoné, à plus haute
température.
\textbf{19.} $10.0/74.0 = \qty{0.135}{mol}$.
\textbf{20.} $1.10 \times 0.135 \times 141.9 = \qty{21.1}{g}$.
\textbf{21.} $0.135 \times 88.0 = \qty{11.9}{g}$.
\textbf{22.} $0.75 \times 11.9 =$ \textbf{\qty{8.9}{g} de MTBE}.
\end{solution}
